\subsection{The determinant passage}

\begin{proposition}\label{sharp:upper-determinant}
Fix \(P_D<\infty\), the passive compact coefficient block and original
window reserves. For independent small \(r_D>0\), \(|\gamma|\), with
\(\alpha_0=r_D^2\), \(|\alpha_1/r_D|\le P_D\), the displacement on
\eqref{deg:det-actual-receiver} has at most one critical point and at most
two isolated zeros, counted with total multiplicity, after a common
coefficient and final multiplier shrink. The raw cap is fixed, positive
and unscaled. The assertion holds for every original window, speed or
height clip and includes saddle ratio one.
\end{proposition}

\begin{proof}
Use the single affine \(F(x)=x-m\Pi(x)\) of Lemma~\ref{lem:closing}.
Lemma~\ref{deg:lem-determinant-preparation} supplies the full stopped
composition, both exact absolute gains and ordinary \(C^3\) divided
closing. On the strict and high sets in \eqref{deg:det-height-cover},
choose \(m\) last so that their complete derivative bounds give
\(F'>1/2\). Thus every critical point lies in the actual stopped range.

There \eqref{sharp:upper-det-absolute} retains every factor of the same
physical derivative:
\[
 \Pi'(x_b+\sigma_+\delta)=\mathcal A_\theta\delta^q e^{R_\theta(\delta)},
 \quad q=\wp-1,\quad
 |R_\theta|+|\delta R_\theta'|\le C\delta^\alpha(1+|\log\delta|).
\]
The quotient of the two gains is a unit, so \(\mathcal A_\theta\)
is bounded above and below; neither absolute gain needs a lower bound.
Let \(C_0\) bound \(|\log\mathcal A_\theta+R_\theta|\).
At a critical point \(m\Pi'=1\), set \(t=-\log\delta>0\). Then
\[
 -qt\ge L:=\log(1/m)-C_0.
\]
For \(L>0\), this requires \(q<0\) and \(|q|\ge L/t\);
in particular \(\wp=1\) produces no critical point. Moreover
\[
 \frac{|\delta R_\theta'|}{|q|}
 \le\frac C L t e^{-\alpha t}(1+t)\le C_1/L.
\]
Choose \(L>2C_1\). The Euler derivative of the logarithm of the full
physical gain satisfies
\(\delta\partial_\delta\log\Pi'<q/2<0\) at every critical point.
Since \(dx/d\delta=\sigma_+>0\),
\begin{equation}\label{sharp:upper-det-critical}
 F'(x)=0\quad\Longrightarrow\quad F''(x)=-m\Pi''(x)>0.
\end{equation}
The actual stopped/strict and low/high cover
\eqref{deg:det-height-cover} retains every clipped and intervening
input. Lemma~\ref{sharp:upper-critical-count} now gives at most
one critical point and two displacement zeros.
\end{proof}

\subsection{The signed trace passage}

\begin{proposition}\label{sharp:upper-trace}
Fix \(0<p_0<p_1<\infty\), the passive compact block and original
window reserves. For \(\alpha_0=r^2a\), \(\alpha_1=rp\),
\(|a|\le a_*\), \(p_0\le|p|\le p_1\), independent small
\(r>0\), \(|\gamma|\), and a final common multiplier shrink,
the displacement on the original receiver
\eqref{deg:trace-actual-receiver} has at most two isolated zeros for
\(p>0\), and at most one for \(p<0\), counted with total multiplicity.
Its positive unscaled raw cap and all window, speed and height clips
are retained. Both opening signs, zero opening and closed trace endpoints
are included.
\end{proposition}

\begin{proof}
Proposition~\ref{deg:prop-stopped-trace} gives \(F'>1/2\) for
\(p<0\) and \(F'=0\Rightarrow F''>0\) for \(p>0\) on the whole
original clipped receiver. Lemma~\ref{sharp:upper-critical-count}
gives the stated multiplicity bounds.
\end{proof}

\begin{corollary}\label{sharp:upper-local-coefficient-ranges}
The nonnegative-opening sector \(\alpha_0\ge0\) has at most
two captured displacement zeros. There is a fixed \(P_*<\infty\)
such that the negative-opening sector
\(\alpha_0=-s^2<0\), \(p_{\rm src}=\alpha_1/s\ge P_*\)
also has at most two. Both assertions concern the whole physical
interval and count multiplicities.
\end{corollary}

\begin{proof}
First choose the signed trace collar of
Proposition~\ref{sharp:upper-trace}, for example with
\(1/2\le|p|\le2\). Choose \(P_D>2/\sqrt{a_*}\).
If \(\alpha_0>0\) and
\(|\alpha_1|\le P_D\sqrt{\alpha_0}\), use
Proposition~\ref{sharp:upper-determinant} with
\(r_D=\sqrt{\alpha_0}\). On the complement put
\(r=|\alpha_1|\). Then
\[
 p=\sgn\alpha_1,\qquad
 a=\alpha_0/\alpha_1^2<P_D^{-2}<a_*.
\]
This is the exact trace scaling of
Lemma~\ref{deg:lem-positive-atlas}; the full physical domains are
retained in Lemma~\ref{deg:lem-nonnegative-domain}. The same choice treats
\(\alpha_0=0\), \(\alpha_1\ne0\), with \(a=0\).
Each alternative bounds the displacement on its whole coefficient fibre.

The double-zero coefficient face is a contraction by
Corollary~\ref{rew:cor-high-derivative-signed} at every positive
height, without introducing a determinant scale.

Finally suppose \(\alpha_0=-s^2<0\) and
\(p_{\rm src}=\alpha_1/s>0\). With \(r=\alpha_1\), the
trace coefficients are exactly
\[
 p=1,\qquad a=-p_{\rm src}^{-2}.
\]
Choose \(P_*>2/\sqrt{a_*}\) before the subsequent compact
positive-source interval is fixed. Then
\(p_{\rm src}\ge P_*\) belongs to
Proposition~\ref{sharp:upper-trace}. Take the maximum of the
whole-fibre bounds across these coefficient alternatives.
\end{proof}
